\section{Background}
\label{sec:background}

\subsection{Word representations}

Let $V$ be the size of the vocabulary. A one-hot representation encodes word $i$ as the $i$-th standard basis vector of $\mathbb{R}^{V}$. These vectors are sparse, as long as the vocabulary, and mutually orthogonal, so they carry no notion of similarity. A word embedding instead assigns each word $w$ a dense vector $\mathbf{v}_w \in \mathbb{R}^{d}$ with $d \ll V$ (here $d = 100$ and $V$ is in the tens of thousands).

Embeddings are learned from the \emph{distributional hypothesis}: words that occur in similar contexts tend to have similar meanings \citep{harris1954distributional}, or in Firth's words, ``you shall know a word by the company it keeps'' \citep{firth1957synopsis}. \textit{King} and \textit{queen} both appear near \textit{crowned}, \textit{reign} and \textit{throne}; a model that has to predict such contexts is pushed to give the two words similar vectors. Neural language models learned embeddings as a by-product of next-word prediction \citep{bengio2003neural}; Word2Vec keeps the idea but strips the model down so that it can be trained on far more text.

\subsection{Word2Vec: CBOW and Skip-gram}

Word2Vec \citep{mikolov2013efficient} is a shallow model with two embedding matrices: an input matrix $W_{\text{in}} \in \mathbb{R}^{V \times d}$ whose row $\mathbf{v}_w$ represents $w$ as a centre word, and an output matrix $W_{\text{out}} \in \mathbb{R}^{V \times d}$ whose row $\mathbf{v}'_w$ represents $w$ as a context word. The score of a (centre $c$, context $o$) pair is the dot product $\mathbf{v}'_o \cdot \mathbf{v}_c$. There are two architectures:
\begin{itemize}
    \item \textbf{CBOW} (continuous bag-of-words) predicts the centre word from the average of its context vectors;
    \item \textbf{Skip-gram} predicts each context word from the centre word.
\end{itemize}
This project uses Skip-gram. Given a corpus $w_1, \dots, w_T$ and a window size $c$, Skip-gram maximises the average log-probability of the context words:
\begin{equation}
    \frac{1}{T} \sum_{t=1}^{T} \sum_{\substack{-c \leq j \leq c \\ j \neq 0}} \log p(w_{t+j} \mid w_t),
    \qquad
    p(o \mid c) = \frac{\exp(\mathbf{v}'_o \cdot \mathbf{v}_c)}{\sum_{w=1}^{V} \exp(\mathbf{v}'_w \cdot \mathbf{v}_c)} .
    \label{eq:skipgram}
\end{equation}
The softmax denominator sums over the whole vocabulary, so every training pair costs $O(V)$ operations. Figure~\ref{fig:skipgram} shows the pairs produced by one centre word.

\begin{figure}[t]
    \centering
    % Skip-gram training pairs for one centre word, with negative sampling (shared by report and slides).
% Requires \usetikzlibrary{positioning,arrows.meta,decorations.pathreplacing}
\begin{tikzpicture}[
    tok/.style={draw=black!35, rounded corners=2pt, minimum height=6.5mm, inner xsep=4pt, font=\small\ttfamily, text=black!45},
    ctx/.style={tok, fill=blue!8, draw=blue!55!black, text=black},
    ctr/.style={tok, fill=orange!22, draw=orange!75!black, text=black, font=\small\ttfamily\bfseries},
    pairbox/.style={draw, rounded corners=3pt, align=left, inner sep=6pt, font=\small},
    >={Stealth[length=2mm]}
]
    \node[tok] (t1) {was};
    \node[ctx, right=2mm of t1] (t2) {crowned};
    \node[ctx, right=2mm of t2] (t3) {the};
    \node[ctr, right=2mm of t3] (t4) {king};
    \node[ctx, right=2mm of t4] (t5) {and};
    \node[ctx, right=2mm of t5] (t6) {queen};
    \node[tok, right=2mm of t6] (t7) {of};
    \node[left=1mm of t1, font=\small, text=black!45] {\dots};
    \node[right=1mm of t7, font=\small, text=black!45] {\dots};

    \draw[->, blue!55!black] ([xshift=-2mm]t4.north) to[bend right=50] (t3.north);
    \draw[->, blue!55!black] ([xshift=-4mm]t4.north) to[bend right=40] (t2.north);
    \draw[->, blue!55!black] ([xshift=2mm]t4.north) to[bend left=50] (t5.north);
    \draw[->, blue!55!black] ([xshift=4mm]t4.north) to[bend left=40] (t6.north);

    \draw[decorate, decoration={brace, mirror, amplitude=4pt}, black!60]
        ([yshift=-1mm]t2.south west) -- ([yshift=-1mm]t6.south east)
        node[midway, below=5pt, font=\footnotesize, text=black!70] (win) {context window of the centre word (here $W = 2$)};

    \node[pairbox, draw=blue!55!black, anchor=north east] (pos) at ([xshift=-3mm, yshift=-4mm]win.south)
        {\textbf{Positive pairs} (label 1)\\[2pt]
         \texttt{(king, crowned)}\\ \texttt{(king, the)}\\ \texttt{(king, and)}\\ \texttt{(king, queen)}};
    \node[pairbox, draw=red!60!black, anchor=north west] (neg) at ([xshift=3mm, yshift=-4mm]win.south)
        {\textbf{Negative pairs} (label 0)\\[2pt]
         \texttt{(king, river)}, \texttt{(king, seven)}, \dots\\[2pt]
         $K$ random words per positive pair,\\
         drawn from $P_n(w) \propto f(w)^{3/4}$};
\end{tikzpicture}

    \caption{Training pairs for the centre word \textit{king}. Each word inside the window forms a positive pair with the centre word; each positive pair is complemented by $K$ negative pairs with words drawn at random from the noise distribution (the negative words shown are illustrative).}
    \label{fig:skipgram}
\end{figure}

\subsection{Negative sampling}
\label{sec:background-ns}

Negative sampling \citep{mikolov2013distributed} replaces the $V$-way softmax with a binary task: tell the observed pair $(c, o)$ apart from $K$ pairs $(c, n_k)$ whose words $n_k$ are drawn from a noise distribution $P_n$. With $\sigma(x) = 1 / (1 + e^{-x})$, the objective for one positive pair is
\begin{equation}
    \log \sigma(\mathbf{v}'_o \cdot \mathbf{v}_c) + \sum_{k=1}^{K} \mathbb{E}_{n_k \sim P_n} \left[ \log \sigma(-\mathbf{v}'_{n_k} \cdot \mathbf{v}_c) \right],
    \qquad
    P_n(w) \propto f(w)^{3/4},
    \label{eq:sgns}
\end{equation}
where $f(w)$ is the frequency of $w$. The first term pulls the centre vector towards the vectors of its true context words; the second pushes it away from random words. Each pair now costs $O(K)$ instead of $O(V)$, with $K$ between 5 and 20 in practice. The exponent $3/4$ samples rare words more often than their raw frequency would. \citet{goldberg2014word2vec} give a detailed derivation of this objective.

The same paper introduces \emph{subsampling} of frequent words: each occurrence of $w$ is discarded with probability $1 - \sqrt{t / f(w)}$ for a threshold $t$ (here $10^{-3}$). Very frequent words such as \textit{the} carry little information about their neighbours, so dropping many of their occurrences speeds up training and gives rarer words more weight.

\subsection{Comparing embeddings}

Embeddings are compared with the cosine similarity
\begin{equation}
    \cos(\mathbf{u}, \mathbf{v}) = \frac{\mathbf{u} \cdot \mathbf{v}}{\lVert \mathbf{u} \rVert \, \lVert \mathbf{v} \rVert} \in [-1, 1] .
\end{equation}
The nearest neighbours of a word are the words with the highest cosine similarity to it. A word analogy ``$a$ is to $b$ as $c$ is to ?'' is answered by vector offsets \citep{mikolov2013linguistic}:
\begin{equation}
    \argmax_{w \notin \{a, b, c\}} \; \cos(\mathbf{v}_w, \; \mathbf{v}_b - \mathbf{v}_a + \mathbf{v}_c),
    \label{eq:analogy}
\end{equation}
so \textit{king} $-$ \textit{man} $+$ \textit{woman} $\approx$ \textit{queen} corresponds to $a = $ \textit{man}, $b = $ \textit{king}, $c = $ \textit{woman}.
